% P26 | 4 Part B summary | What Part B shows about the European call
% guide: dev_notes/paragraph_guides/P26_what_part_b_shows_about_the_european_call.txt
For a single European call on one stock, a Leisen--Reimer tree or a smoothed and extrapolated CRR tree is orders of magnitude more accurate than Monte Carlo for the same run time.
The work--precision comparison in Section~\ref{sec:compare} puts numbers on this.
Monte Carlo wins when the payoff depends on the whole path or on many underlyings, because its error rate does not depend on the dimension, while the cost of a lattice grows exponentially with it.
Variance reduction closes part of the gap: a control variate alone was worth a factor of \VRcv\ in variance here.
The three routes to the Greeks agree, so the choice among them is a matter of cost and variance, not of correctness.
The hedging error of a short call falls only like $n^{-1/2}$ with the number of rebalances, and an error of $10$ volatility points in the hedge volatility costs more than any reasonable change of frequency.
Early exercise needs a decision at every date based on the value of waiting, which the tree computes naturally and Monte Carlo can approximate only with a regression, as Part C shows.
